\section{Модели взаимодействия с обучающимися}
\thispagestyle{plain}

% TODO: ~16 страниц --- ключевая глава по теме работы.

\subsection{Ролевая модель пользователей}
% Преподаватель, студент, администратор. Атрибуты, права доступа,
% диаграмма прецедентов (use case diagram).
% Мультитенантность: несколько преподавателей не пересекаются
% по своим курсам.

\subsection{Модель курса и банка вопросов}
% Курс (course) --- логическая единица, связанная с дисциплиной.
% Банк вопросов (question bank) --- множество вопросов,
% сгруппированных по темам.
% Типы вопросов: с одиночным выбором, с множественным выбором,
% открытый ответ, рейтинг, Q\&A.
% Метаданные вопроса: тема, сложность, теги, источник
% (импортирован/создан вручную).

\subsection{Модель интерактивной сессии}
% Сессия (room) --- развёрнутая в реальном времени викторина или опрос.
% Состояния: создана, ожидание, активна (с конкретным вопросом),
% завершена.
% Уникальный код доступа, протокол подключения студентов,
% отказоустойчивость подключения.

\subsection{Модель игровой прогрессии (геймификация)}
% Очки (XP), уровни, таблица лидеров.
% Бейджи (achievements): за серии правильных ответов,
% завершение тем, активное участие.
% Серии (streaks): ежедневная активность.
% Формула начисления XP с учётом скорости ответа и сложности вопроса.

\subsection{Модель интервального повторения (SM\textendash 2)}
% Алгоритм SuperMemo SM-2 (Wozniak, 1990).
% Параметры карточки: easiness factor (E), interval (I), repetitions (n).
% Формулы обновления при разных оценках качества ответа.
% Применение к закреплению знаний теории надёжности на курсе Волкова Д.А.

\subsection{Модель оценки надёжности усвоения знаний}
% Альфа Кронбаха (Cronbach $\alpha$) --- внутренняя согласованность
% тестового инструмента.
% Элементы IRT (Item Response Theory): двухпараметрическая модель Раша.
% Метрики удержания (retention) при повторном тестировании.

\subsection{UML-диаграммы}
% Диаграмма классов основных сущностей.
% Диаграмма последовательностей для сценария <<студент отвечает на вопрос>>.
% Диаграмма состояний интерактивной сессии.

\textbf{Выводы по главе 2.}
% Какие модели предложены, чем они отличаются от существующих,
% как они связаны с теорией надёжности (методический уровень).

\newpage
